\chapter{Dynamic Programming}\index{dynamic programming}

Type ``libary'' into your phone, and it suggests ``library.'' Run \verb|git diff| after editing your
essay, and Git shows exactly which lines you added and which you deleted. Biologists line up the DNA of
a human and a chimpanzee to see which genes they share. All three are the same problem: comparing two
sequences to find what they have in common.

The obvious way to solve it, trying every possibility, would take longer than the age of the universe
for an essay of any real length. The trick that makes it fast is called \newterm{dynamic programming}:
break a problem into smaller pieces that overlap, solve each piece \emph{once}, and write the answers in
a table so you never solve it again. In this chapter, you will use it to find the longest common
subsequence of two strings, and then turn that into a working \verb|diff| tool.

\section{Remember Fibonacci?}
In the Recursion chapter, you saw that the obvious recursive Fibonacci function is painfully slow,
because \pyfunction{fib(n - 1)} and \pyfunction{fib(n - 2)} recompute the same smaller values over and
over. You fixed it with \newterm{memoization}: save each answer in a dictionary the first time you
compute it, and look it up after that.

There is a second way to get the same speedup. Instead of starting at the top and recursing down, start
at the bottom and fill in a table, smallest values first:

\begin{minted}{python}
def fib(n):
    table = [0, 1]
    for k in range(2, n + 1):
        table.append(table[k - 1] + table[k - 2])
    return table[n]

print(fib(80))   # 23416728348467685
\end{minted}

By the time the loop needs \verb|table[k - 1]| and \verb|table[k - 2]|, they are already filled in. No
recursion, no repeated work: just 79 additions.

These are the two styles of dynamic programming. Memoization is \newterm{top-down}: recurse as usual, but
remember answers. Filling a table is \newterm{bottom-up}: solve the small problems first, in an order that
guarantees each answer is ready when it is needed.

Dynamic programming works whenever a problem has two features:
\begin{itemize}
\item \textbf{Optimal substructure}: the answer to the big problem can be built from answers to smaller
versions of the same problem.
\item \textbf{Overlapping subproblems}: the same smaller problems come up again and again. (If they did not
overlap, plain recursion would already be fast.)
\end{itemize}

\section{Longest Common Subsequence}\index{longest common subsequence}
A \newterm{subsequence} of a string is what is left after you cross out some of its letters, without
changing the order of the ones that remain. Crossing out letters from STUDENT can give DENT, SUN, or
STENT, but not NUTS, because that changes the order.

A \newterm{common subsequence} of two strings is a subsequence of both. The \newterm{longest common
subsequence} (LCS) is the longest one there is. For STUDENT and DENTIST, it is DENT:
\begin{center}
S T U \textbf{D E N T} \qquad\qquad \textbf{D E N T} I S T
\end{center}

How could you find it? You could list every subsequence of the first string and check whether each one is
a subsequence of the second. But each letter is either kept or crossed out, so a string of $n$ letters has
$2^n$ subsequences. For a 20-letter string, that is over a million. For a 1,000-line essay compared line by
line, $2^{1000}$ is a number with 302 digits. Brute force is hopeless.

\subsection{Thinking Recursively}
Look at just the \emph{last} letters of the two strings. There are two cases.

\textbf{The last letters match.} Then that letter can end the LCS. Chop it off both strings, find the LCS
of what is left, and add 1. For example, the LCS of SHOE and HOSE is 1 more than the LCS of SHO and HOS,
because both end in E.

\textbf{The last letters do not match.} Then at least one of them is not in the LCS. We do not know which,
so try both: drop the last letter of the first string, or drop the last letter of the second string, and
keep whichever gives the longer answer.

If either string is empty, the LCS has length 0. That is the base case. Here it is in Python, in a file
called \filename{lcs\_recursive.py}:

\begin{minted}{python}
calls = 0

def lcs_length(a, b):
    global calls
    calls += 1
    if a == "" or b == "":
        return 0
    if a[-1] == b[-1]:
        return lcs_length(a[:-1], b[:-1]) + 1
    return max(lcs_length(a[:-1], b), lcs_length(a, b[:-1]))

print(lcs_length("SHOE", "HOSE"), "after", calls, "calls")
calls = 0
print(lcs_length("DYNAMICPROG", "PROGRAMMING"), "after", calls, "calls")
\end{minted}

When you run it, you should see:
\begin{verbatim}
3 after 16 calls
4 after 212591 calls
\end{verbatim}

The answers are right, but two 11-letter words took over 200,000 calls. When the last letters do not match,
the function calls itself \emph{twice}, and those two calls ask many of the same smaller questions. Just
like Fibonacci, the work grows exponentially. And just like Fibonacci, the fix is to stop recomputing.

\subsection{Filling the Table}
Every smaller question has the same form: what is the LCS of the first $i$ letters of $a$ and the first $j$
letters of $b$? Call the answer $L[i][j]$. The recursive reasoning becomes a rule for filling a table:

\begin{mdframed}[style = important, frametitle = {The LCS Recurrence}]
$$L[i][j] = \begin{cases}
0 & \text{if } i = 0 \text{ or } j = 0\\
L[i-1][j-1] + 1 & \text{if the } i\text{th letter of } a \text{ matches the } j\text{th letter of } b\\
\max\big(L[i-1][j],\ L[i][j-1]\big) & \text{otherwise}
\end{cases}$$
\end{mdframed}

Each cell needs only the cell above it, the cell to its left, and the cell diagonally up and to the left. So
if you fill the table row by row, left to right, every cell you need is already done.

\textbf{Example}: Fill in the table for $a = $ SHOE (down the side) and $b = $ HOSE (across the top).

\textbf{Solution}: The first row and column are all 0, because one of the strings is empty. Then fill in each
row. In row S, the only match is with the S in HOSE, so that cell is $0 + 1 = 1$, and the 1 is carried to
the right. In row H, the H matches the first column, giving 1. In row O, the O matches the second column:
the diagonal cell is 1, so this cell is 2. In row E, the E matches the last column, and the diagonal cell is
2, so the last cell is 3.

\begin{center}
\begin{tabular}{c|ccccc}
 & & H & O & S & E \\
\hline
 & 0 & 0 & 0 & 0 & 0 \\
S & 0 & 0 & 0 & 1 & 1 \\
H & 0 & \textbf{1} & 1 & 1 & 1 \\
O & 0 & 1 & \textbf{2} & 2 & 2 \\
E & 0 & 1 & 2 & 2 & \textbf{3} \\
\end{tabular}
\end{center}

The bottom-right cell says the LCS has length 3.

To find the letters themselves, walk backward from the bottom-right corner. If the letters for this cell
match, that letter is part of the LCS: record it and step diagonally. If not, step toward whichever neighbor
(up or left) holds the larger number. The bold cells are where the walk finds matches: H, O, and E, so the
LCS is HOE. This is the same idea as following the \verb|prev| links backward in Dijkstra's algorithm.

\subsection{The Code}
Create a file called \filename{lcs.py}:

\begin{minted}{python}
def lcs_table(a, b):
    # table[i][j] = length of the LCS of a[:i] and b[:j]
    table = [[0] * (len(b) + 1) for i in range(len(a) + 1)]
    for i in range(1, len(a) + 1):
        for j in range(1, len(b) + 1):
            if a[i - 1] == b[j - 1]:
                table[i][j] = table[i - 1][j - 1] + 1
            else:
                table[i][j] = max(table[i - 1][j], table[i][j - 1])
    return table

def lcs(a, b):
    table = lcs_table(a, b)
    # Walk back from the bottom-right corner to recover the letters
    i, j = len(a), len(b)
    letters = []
    while i > 0 and j > 0:
        if a[i - 1] == b[j - 1]:
            letters.append(a[i - 1])
            i -= 1
            j -= 1
        elif table[i - 1][j] >= table[i][j - 1]:
            i -= 1
        else:
            j -= 1
    letters.reverse()
    return "".join(letters)

if __name__ == "__main__":
    print(lcs("SHOE", "HOSE"))
    print(lcs("STUDENT", "DENTIST"))
    print(lcs("DYNAMICPROG", "PROGRAMMING"))
\end{minted}

The $i$th letter of \verb|a| is \verb|a[i - 1]|, because Python counts from 0 while the table's row 0 stands
for the empty string. The \verb|if __name__ == "__main__":| line, which you saw in the Testing and Debugging
chapter, means ``run these lines only when this file is run directly.'' That lets another program import
\pyfunction{lcs\_table} without printing anything.

When you run it, you should see:
\begin{verbatim}
HOE
DENT
AMIG
\end{verbatim}

The pair that needed 212,591 recursive calls now takes a $12 \times 12$ table: 144 cells. (There can be more
than one longest common subsequence; \verb|AMIG| is one of several of length 4.)

How fast is it in general? The table has $(n + 1)(m + 1)$ cells, and each takes a few steps, so the time
grows like $nm$. On a laptop, comparing two random DNA strands of 2,000 letters each (4 million cells) takes
a fraction of a second. Exponential has become polynomial.

\section{Building a Diff}\index{diff}
In the Version Control with Git chapter, Git kept every version of your files. Run \verb|git diff| in any
repository after editing a file, and Git shows what changed since your last commit. It does this by finding the longest common subsequence of the two versions, using whole \emph{lines} as the
``letters.'' Lines in the LCS are unchanged. Lines that are only in the old version were deleted, and lines
that are only in the new version were added.

Since Python compares lists of lines with \verb|==| just as it compares letters, \pyfunction{lcs\_table}
works on lists without any changes. Create \filename{diff.py} in the same folder:

\begin{minted}{python}
from lcs import lcs_table

def diff(old, new):
    table = lcs_table(old, new)
    i, j = len(old), len(new)
    lines = []
    while i > 0 or j > 0:
        if i > 0 and j > 0 and old[i - 1] == new[j - 1]:
            lines.append("  " + old[i - 1])     # in both: unchanged
            i -= 1
            j -= 1
        elif j > 0 and (i == 0 or table[i][j - 1] >= table[i - 1][j]):
            lines.append("+ " + new[j - 1])     # only in new: added
            j -= 1
        else:
            lines.append("- " + old[i - 1])     # only in old: deleted
            i -= 1
    lines.reverse()
    return lines

monday = ["Buy poster board", "Finish lab report", "Email Ms. Rivera",
          "Practice trumpet", "Return library book"]
tuesday = ["Buy poster board", "Finish lab report", "Practice trumpet",
           "Study for chem quiz", "Return library book"]

for line in diff(monday, tuesday):
    print(line)
\end{minted}

The walk is the same as in \pyfunction{lcs}, except that it keeps going until \emph{both} lists are used up,
and every step it takes becomes a line of output. When you run it, you should see:
\begin{verbatim}
  Buy poster board
  Finish lab report
- Email Ms. Rivera
  Practice trumpet
+ Study for chem quiz
  Return library book
\end{verbatim}

That is the same format Git uses. Real diff tools add a few refinements, but the heart of every one of them is
this table.

\section{The Dynamic Programming Recipe}
Most dynamic programming solutions are built in the same five steps:
\begin{enumerate}
\item \textbf{Define the subproblem.} Say in words what one cell of the table means. (``$L[i][j]$ is the length
of the LCS of the first $i$ letters of $a$ and the first $j$ letters of $b$.'')
\item \textbf{Write the recurrence.} Express each cell in terms of smaller cells.
\item \textbf{Set the base cases.} Fill in the cells that need no smaller ones, such as the empty strings.
\item \textbf{Choose a filling order.} Make sure every cell's ingredients are ready before you need them.
\item \textbf{Read off the answer,} and walk back through the table if you need more than a number.
\end{enumerate}

You have met dynamic programming before without its name. The Bellman-Ford algorithm in the Finding Shortest
Paths chapter is dynamic programming: after round $k$, each distance is the shortest path using at most $k$
edges, built from the round before. The Viterbi algorithm, mentioned in the Introduction to Markov Chains chapter, is
dynamic programming too.

The name, by the way, is not very descriptive. Richard Bellman (the same Bellman) chose it in the 1950s, when the
government officials who funded his work distrusted mathematical research. He later wrote that he picked
``dynamic'' partly because it sounded impressive and could not be used as an insult.
``Programming'' meant planning, as in a TV program schedule, not writing code.

\begin{Exercise}[title={By Hand}, label=dp_hand]
Fill in the LCS table for CRATE (down the side) and TRACE (across the top). What is the length of the LCS?
Walk back through your table to find one.
\end{Exercise}
\begin{Answer}[ref=dp_hand]
\begin{tabular}{c|cccccc}
 & & T & R & A & C & E \\
\hline
 & 0 & 0 & 0 & 0 & 0 & 0 \\
C & 0 & 0 & 0 & 0 & 1 & 1 \\
R & 0 & 0 & 1 & 1 & 1 & 1 \\
A & 0 & 0 & 1 & 2 & 2 & 2 \\
T & 0 & 1 & 1 & 2 & 2 & 2 \\
E & 0 & 1 & 1 & 2 & 2 & 3 \\
\end{tabular}

The length is 3. Walking back with the rules from the chapter gives RAE, and in fact RAE is the only common
subsequence of length 3. (RCE does not work, because C comes before R in CRATE.)
\end{Answer}

\begin{Exercise}[title={Climbing Stairs}, label=dp_stairs]
You can climb a staircase 1 or 2 steps at a time. How many different ways are there to climb 10 steps? (Hint:
your last move was either 1 step or 2 steps. Let $W[n]$ be the number of ways to climb $n$ steps.)
\end{Exercise}
\begin{Answer}[ref=dp_stairs]
The last move came from step $n - 1$ or step $n - 2$, so $W[n] = W[n-1] + W[n-2]$, with $W[0] = 1$ (do nothing)
and $W[1] = 1$. Filling the table: 1, 1, 2, 3, 5, 8, 13, 21, 34, 55, 89. There are 89 ways. These are the
Fibonacci numbers again.
\end{Answer}

\begin{Exercise}[title={Making Change}, label=dp_coins]
You need to make 30 cents using pennies (1), dimes (10), and quarters (25), with as few coins as possible.
\begin{enumerate}
\item The \emph{greedy} method always takes the biggest coin that fits. How many coins does it use?
\item What is the true minimum?
\item Write a function \pyfunction{fewest\_coins(amount, coins)} that uses a table \verb|best|, where
\verb|best[t]| is the fewest coins that make \verb|t| cents.
\end{enumerate}
\end{Exercise}
\begin{Answer}[ref=dp_coins]
\begin{enumerate}
\item Greedy takes a quarter, then 5 pennies: 6 coins.
\item Three dimes: 3 coins. Taking the biggest coin first is not always best.
\item For each total, try ending with each coin that fits, and keep the best:
\begin{minted}{python}
import math

def fewest_coins(amount, coins):
    best = [0] + [math.inf] * amount
    for total in range(1, amount + 1):
        for coin in coins:
            if coin <= total and best[total - coin] + 1 < best[total]:
                best[total] = best[total - coin] + 1
    return best[amount]

print(fewest_coins(30, [1, 10, 25]))   # 3
\end{minted}
\end{enumerate}
\end{Answer}

\begin{Exercise}[title={Spell Check}, label=dp_spell]
Suppose the only way to fix a typo is to delete letters and insert letters. Explain why the fewest changes
needed to turn word $a$ into word $b$ is $\text{len}(a) + \text{len}(b) - 2 \cdot \text{LCS}$. Use it to find
the number of changes from ``libary'' to ``library'' and from ``recieve'' to ``receive.''
\end{Exercise}
\begin{Answer}[ref=dp_spell]
Keep the letters of the LCS and change everything else: delete the $\text{len}(a) - \text{LCS}$ letters of $a$
that are not in it, then insert the $\text{len}(b) - \text{LCS}$ letters of $b$ that are missing. That adds up
to $\text{len}(a) + \text{len}(b) - 2 \cdot \text{LCS}$, and no plan can do better, because any letters you
never touch form a common subsequence.

``libary'' and ``library'' have an LCS of length 6, so the count is $6 + 7 - 12 = 1$: insert an r.
``recieve'' and ``receive'' have an LCS of length 6, so the count is $7 + 7 - 12 = 2$: delete the i and insert
it after the e. A spell checker can suggest the dictionary words with the smallest counts.
\end{Answer}

\begin{Exercise}[title={How Big Is the Table?}, label=dp_size]
You compare two essays of 1,000 lines each, line by line. How many cells does the LCS table have? How many
subsequences does brute force have to try?
\end{Exercise}
\begin{Answer}[ref=dp_size]
The table has $1{,}001 \times 1{,}001 = 1{,}002{,}001$ cells, about a million, which a laptop fills in well under
a second. Brute force would have to try $2^{1000}$ subsequences of one essay, a number with 302 digits.
\end{Answer}

\begin{Exercise}[title={Substring or Subsequence?}, label=dp_substring]
A \emph{substring} must be made of letters that are next to each other, like DENT in STUDENT. What is the
longest common substring of SNOWBOARD and SKATEBOARD? What is their longest common subsequence?
\end{Exercise}
\begin{Answer}[ref=dp_substring]
The longest common substring is BOARD (5 letters, side by side in both). The longest common subsequence is
SBOARD (6 letters): the S at the start of each word can join BOARD, because a subsequence may skip letters.
\end{Answer}
